\section{Teaching Experience}

\subsection{At Johns Hopkins University}
\begin{cvlist}
  \cvitem{2021-2026}{\entrytitle{Instructor, Machine Learning for Medical Applications (EN.520.439/659)}
\begin{itemize}
    \item Developed a new senior undergraduate/graduate level course.
    \item In this course, students actively learn the basic principles of artificial intelligence and machine learning techniques applied to medical applications. Throughout the course, students also discover new concepts about different types of bio-signals such as electroencephalograms, medical imaging, and their associated processing methodologies. The primary objective is to give students the tools they need to be able to develop new ideas relating to biomedical environments or other fields. At the end of the course, students will apply their newly acquired knowledge to complete a cumulative final project.
\end{itemize}}
  \cvitem{2021-2022}{\entrytitle{Instructor, AI in Medicine – Reading Group (EN.520.520/820)}
\begin{itemize}
    \item Developed a new graduate level course.
    \item In this course, the students learn how to analyze and understand scientific papers related to artificial intelligence in medical applications. The students also learn how to review the studies, perform critical reviews, and present and discuss their analysis in front of the class.
\end{itemize}}
  \cvitem{2021-2025}{\entrytitle{Instructor, Introduction to Human Language Technology   
  (EN.601.467/667)}
\begin{itemize}
    \item Instructor of the module dealing with ethical problems in artificial intelligence.
    \item Instructor of the module Speech Production (since 2022)
\end{itemize}}
  \cvitem{2019}{\entrytitle{Instructor,  HEART: Machine learning for biomedical applications (EN.500.111-14/23)}
\begin{itemize}
    \item Developed a new course for first-year students.
    \item In this course, students intuitively and actively learn the basic principles of machine learning techniques such as K-Nearest Neighbors, Gaussian Mixture Models and Neural Networks. We will also discuss about their potential use in real-life applications but specifically focusing on medical applications. The primary objective is to give students the basic tools they need to be able to develop new ideas relating to biomedical environments.
\end{itemize}}
  \cvitem{2018}{\entrytitle{Teaching assistant, Machine Learning for Signal Processing (EN.520.412/612)}}
  \cvitem{2018}{\entrytitle{Teaching assistant, Computational Modelling for Electrical and Computer Engineering (EN.520.123)}}
\end{cvlist}

\subsection{At Universidad Politécnica de Madrid}
\begin{cvlist}
  \cvitem{2016}{\entrytitle{Instructor, Signals and systems}
\begin{itemize}
    \item Course for senior undergraduate students
\end{itemize}}
  \cvitem{2016}{\entrytitle{Instructor,  Circuit Analysis Laboratory}
\begin{itemize}
    \item Course for first-year undergraduate students
\end{itemize}}
  \cvitem{2010 – 2011}{\entrytitle{Instructor, Acoustic Instrumentation}
\begin{itemize}
    \item Developed the laboratory section of a new course for graduate students as part of the UPM Acoustical Engineering for Construction and Environment Master’s degree.
    \item In this course, students acquire an advanced knowledge of acoustic instrumentation to be used in environmental and architectural measurements.
\end{itemize}}
  \cvitem{2009 – 2011}{\entrytitle{Instructor, Audiofrequency Systems I and II}
\begin{itemize}
    \item Course for senior undergraduate students
\end{itemize}}
\end{cvlist}
